\section{Preliminaries}
\label{sec:preliminaries}

All logarithms are in base $2$.  For integers $n \ge 1$ we write
$[n] = \set{1, 2, \ldots, n}$ and for a set $U$ we denote the power
set of $U$ by $2^{U}$. Further, for a set $V \subseteq U$ we let
$\overline V$ be the complement of $V$ with respect to $U$, that is,
$\overline V = U \setminus V$. We write
$\binom{[n]}{\ell} \subseteq \set{0,1}^n$ for the set of binary
strings with Hamming weight $\ell$. For a string
$\alpha \in \set{0,1}^n$ we let
$|\alpha| = \sum_{i \in [n]} \alpha_i$.

We sometimes want to supress dependencies on constants and write
$f(n, \eps) \in O_\eps\big(g(n, \eps)\big)$, respectively
$f(n, \eps) \in \Omega_\eps\big(g(n, \eps)\big)$, to mean that there
exists a function $c(\eps) > 0$ such that there is an $n_0$ and for all
$n \ge n_0$ it holds that $f(n, \eps) \le c(\eps) \cdot g(n, \eps)$,
respectively $f(n, \eps) \ge c(\eps)\cdot g(n, \eps)$.

\begin{definition}\label{def:explicit}
  A sequence of bipartite graphs $\set{G_n=(U_n, V_n, E_n)}_{n\in \N}$
  with $\deg(u) =d$ for all $u \in U_n$ is \emph{explicit} if there is
  an algorithm that given $(n, u, j)$, where $n \in \N, u \in U_n$ and
  $j\in [d]$, computes the $j$th neighbor of vertex $u$ in the graph
  $G_n$ in time $\poly(\log n + \log |U| + \log d)$.
\end{definition}

From now on it is understood that whenever we talk about an explicit
graph we actually mean to say that there is a sequence of explicit
graphs with above properties.

For a graph $G = (V, E)$ and $W \subseteq V$ we denote by
\[
  N(W) = \setdescr{u \in V \setminus W}{(w,u) \in E \text{ for some } w
    \in W}
\]
the set of neighbors of $W$.

\begin{definition}
  A bipartite graph $G = (U, V, E)$ is an $(r, d, c)$-expander if every
  vertex $u \in U$ has degree $\deg(u) = d$ and every set
  $W \subseteq U$ of size $|W| \le r$ satisfies $|N(W)| \ge c \cdot |W|$.
\end{definition}

A key ingredient in our proofs is the following result on the
existence of strong explicit expanders.

\begin{theorem}[\cite{GUV09Unbalanced}]\label{thm:expander}
  For all constants $\gamma > 0$,
  every $M \in \N$, $r \le M$,
  and $\eps > 0$,
  there is an
  $N \le d^2 \cdot r^{1 + \gamma}$ and an
  explicit
  $(r, d, (1- \eps)d)$-expander
  $G=(U, V, E)$,
  with $|U| = M$,
  $|V| = N$, and
  $d = O\big( ((\log M)(\log r) / \eps )^{1+1/\gamma}\big)$.
\end{theorem}

\kr{Igor pointed this out:
  \url{https://eccc.weizmann.ac.il/report/2018/159/}. Quickly checking
  the parameters I believe it does not give anything better. Need to
  check more carefully}

For our purposes it is more relevant to compute the neighbor relation
$\neigh(u,v)$ indicating whether $(u, v) \in E$ rather than the
neighbor function as in \cref{def:explicit}, but this is an immediate
consequence of being able to compute the neighbor function.

\begin{claim}
  \label{claim:explicit neigh}
  If $G = (U, V, E)$ is explicit then the neighbor relation
  $\neigh: U \times V \rightarrow \{0,1\}$ is computable by a circuit
  of size
  $d \cdot \big( \poly(\log n + \log |U| + \log d) + 2 \log |V| + 1
  \big)$.
\end{claim}

% \pa{Added this fact, check that I was not too optimistic in claimed
  % size bound. Maybe add proof here.} \kr{checked.}

% Recall that this graph is explicit and hence there are
% circuits to compute the neighbor function
% $\Gamma: U \times [k] \rightarrow Z$.

% Each circuit $C_1, \ldots, C_z$ consists of $k$ circuits $\calD =
% \set{D_1, \ldots, D_k}$ that compute $\Gamma$ with the second input
% fixed to their corresponding index. On top of these circuits from
% $\calD$ we add a small circuit  $E_i$, which computes the following.
% Each circuit $D \in \calD$ outputs
% some index $j \in [z]$. The circuit $E_i$ outputs $1$ if and only if
% one of the $D \in \calD$ outputs the index $i$. The circuit $E_i$ is
% fairly small: for each circuit $D$ we require $2 \log z$ gates to
% check whether $D$ output $i$, and an additional $k$ gates to \emph{or}
% these outputs together. So in total each $E_i$ consists of
% $k (2\log z + 1)$ many gates.
% As the graph $G$ is explicit, we know that there is a circuit of size
% $\poly(n + \log k)$ computing $\Gamma$ and hence the size of all
% circuits in $\calD$ is bounded by $c_\alpha \cdot k \cdot n^c$, for
% some constant $c$, and $c_\alpha$ depending on $\alpha$. Hence we
% conclude that
% \begin{align}\label{eq:bound-a}
%  a = z \cdot \big(k (2\log z + 1) + c_\alpha \cdot k \cdot n^c\big) \eqperiod
% \end{align}


A slice function is a Boolean function $f$ such that there is a
$\ell \in [n]$ with $f(\alpha) = 0$ whenever $|\alpha| < \ell$, and
$f(\alpha) = 1$ whenever $|\alpha| > \ell$.  Note that all slice
functions are monotone.

The circuit complexity $\csize(f)$ of a Boolean function $f$ is the
size of the smallest circuit over the basis $\vee, \wedge$, and
$\lnot$ (with fan-in $2$). Similarly the monotone circuit
complexity $\cmonsize(f)$ of a monotone Boolean function $f$
is the size of the smallest circuit over the basis $\vee$, and
$\wedge$. We have the following useful inequality between these
measures.

\begin{lemma}[\cite{B82}]\label{lem:slice}
  If $g$ is any slice function on $n$ bits, then
  $\cmonsize(g) \le 2 \csize(g) + O(n^2 \log n)$.
\end{lemma}

Finally we also rely on the following simple claim.

\begin{claim}
  \label{claim:0 mod bool}
  Let $p: \R^n \rightarrow \R$ be a degree $d$ polynomial such that
  $p(x) = 0$ for all $x \in \{0,1\}^n$.  Then $p$ can be written as
  \[
    p(x) = \sum_{i \in [n]} q_i(x) \cdot (x_i^2-x_i)
  \]
  where each term in the sum has degree at most $d$.
\end{claim}

\begin{proof}[Proof sketch]
  We take the polynomial $p$ and multilinearize it, using the
  appropriate polynomial $x_i^2 - x_i$. Eventually we are left with a
  sum of polynomials of the form $q_i(x)\cdot(x_i^2 - x_i)$ and a
  multilinear polynomial $\tilde p(x)$ which is $0$ on all Boolean
  inputs.  As multilinear polynomials are a basis for Boolean
  functions this implies that $\tilde p(x)$ is equal to the $0$
  polynomial and hence the claim follows.
\end{proof}

\subsection{Sum of Squares}
\label{sec:known-SoS-results}

Let $\calP = \set{p_1 = 0, \ldots, p_m = 0}$ be a system of polynomial
equations over the set of variables
$X = \set{x_1, \ldots, x_n, \bar{x}_1, \ldots, \bar{x}_n}$.  Each
$p_i$ is called an \emph{axiom}, and throughout the paper we always
assume that $\calP$ includes all axioms $x_i^2 - x_i$ and
$\bar{x}_i^2 - \bar{x}_i$, ensuring that the variables are Boolean, as
well as the axioms $1 - x_i - \bar{x}_i$, making sure that the ``bar''
variables are in fact the negation of the ``non-bar'' variables.

%

\begin{definition}[Sum-of-Squares]
  \emph{Sum-of-Squares (SoS)} is a static, semi-algebraic proof
  system. An SoS proof of $f \ge 0$ from $\calP$ is a sequence of
  polynomials $\pi = (t_1, \ldots, t_m; s_1, \ldots, s_a)$ such that
  \begin{align*}
    \sum_{i \in [m]} t_i p_i +
    \sum_{i \in [a]} s^2_i = f \eqperiod %\label{eq:sos-def}
  \end{align*}
  The \emph{degree} of a proof $\pi$ is
  $\Deg(\pi) = \max \set{ \max_{i \in [m]} \deg(t_i) + \deg(p_i),
    \max_{i \in [a]} 2 \deg(s_i) }$, an \emph{SoS refutation of
    $\calP$} is an SoS proof of $-1 \ge 0$ from $\calP$, and the
  \emph{SoS degree to refute $\calP$} is the minimum degree of any SoS
  refutation of $\calP$: if we let $\pi$ range over all SoS
  refutations of $\calP$, we can write
  $\RefuteDeg{SoS}{\calP} = \min_\pi \Deg(\pi)$.  The \emph{size of an
    SoS refutation $\pi$}, $\Size(\pi)$, is the sum of the number of
  monomials in each polynomial in $\pi$ and the \emph{size of refuting
    $\calP$} is the minimum size over all refutations
  $\RefuteSize{SoS}{\calP} = \min_{\pi} \Size(\pi)$.
\end{definition}


Let us recall some well-known results about SoS. Given a bipartite
graph $G = (U, V, E)$, and $b \in \set{0,1}^{|U|}$ we denote by
$\Phi(G,b)$ the following XOR-CSP instance defined over $G$: for each
$v \in V$ there is a Boolean variable $x_v$, and for every vertex
$u \in U$ there is a constraint $\oplus_{v \in N(u)} x_v = b_u$.
We encode this in the obvious way as a system of polynomial equations: 
\begin{align*}
  \big\{
  \prod_{v \in N(u)} (1 - 2 \cdot x_v) = 1 - 2 \cdot b_u
  \mid u \in U
  \big\} \eqcomma
\end{align*}
along with the Boolean axioms and the negation axioms for the $x$
variables. The first theorem we need to recall is the classic lower
bounds for XOR-CSPs by Grigoriev.

\begin{theorem}[\cite{gri01xor}]\label{thm:sos-xor}
  For $n \in \N$, all $k = k(n)$ and $r = r(n)$ the following holds.
  Let $G = (U, V, E)$ be an $(r, k, 2)$-expander with $|V| = n$. Then
  for every $b \in \set{0,1}^{|U|}$ SoS requires degree $\Omega(r)$ to
  refute the claim that there is a satisfying assignment to
  $\Phi(G, b)$.
\end{theorem}

We also need to recall the size-degree tradeoff by Atserias and
Hakoniemi.

\begin{theorem}[\cite{AH18SosTradeoff}]\label{thm:tradeoff}
  Let $\calP$ be a system of polynomial equations over $n$ Boolean
  variables and degree at most $k$. If $d$ is the minimum degree
  SoS requires to refute $\calP$, then the minimum size of an SoS
  refutation of $\calP$ is at least $\exp(\Omega((d-k)^2/n))$.
\end{theorem}


\subsection{Restrictions}
\label{sec:restrictions}

Let $\calP = \set{p_1 = 0, \ldots, p_m = 0}$ be a system of polynomial
equations over the set of Boolean variables
$X = \set{x_1, \ldots, x_n, \bar{x}_1, \ldots, \bar{x}_n}$.
%
For a map
$\rho: \set{x_1, \ldots, x_n} \to \set{0,1, x_1, \ldots, x_n,
  \bar{x}_1, \ldots, \bar{x}_n}$ denote by $\restrict{\calP}{\rho}$
the system of polynomial equations $\calP$ restricted by $\rho$, i.e.,
\begin{align*}
\restrict{\calP}{\rho} = \set{ &p_1(\rho(x_1), \ldots, \rho(x_n)) = 0,\\
  &p_2(\rho(x_1), \ldots, \rho(x_n)) = 0,\\ &\vdots\\ &p_m(\rho(x_1),
  \ldots, \rho(x_n)) = 0 }\eqcomma
\end{align*}
where it is understood that $\rho(\bar x_i) = \overline{\rho(x_i)}$,
with the convention $\bar{\bar{x}}_i = x_i$, $\bar 0 = 1$ and vice
versa. Throughout the paper all our restrictions set the bar variables
to the negation of the non-bar variables. As such it makes sense to
treat the pair of variables $(x_i, \bar x_i)$ as one variable and we
say that $\calP$ has $n$ \emph{unset} variables.

\begin{definition}[Variable Substitution]
  We say that a system of polynomial equations $\calP'$ is a
  \emph{variable substitution of $\calP$} if there is a map
  $\rho: \set{x_1, \ldots, x_n} \to \set{0,1, x_1, \ldots, x_n,
    \bar{x}_1, \ldots, \bar{x}_n}$ such that
  $\calP' = \restrict{\calP}{\rho}$, where we ignore polynomial
  equations of the form $0=0$.
\end{definition}

The following well-known lemma states that a system of polynomial
equations $\calP$ is at least as hard as any of its variable substitutions.
\begin{lemma}\label{lem:subformula}
  Let $\calP, \calP'$ be systems of polynomial equations such that
  $\calP' \subformula \calP$. Then,
  \begin{enumerate}[label=$(\roman*)$]
  \item $\RefuteDeg{SoS}{\calP} \ge \RefuteDeg{SoS}{\calP'}$, and
  \item $\RefuteSize{SoS}{\calP} \ge \RefuteSize{SoS}{\calP'}$.
  \end{enumerate}
\end{lemma}
The lemma is easy to verify by considering an SoS refutation of
$\calP$ and hitting it with the appropriate variable substitution. The
restricted proof is now a refutation of $\calP'$ and it can be seen
that the degree/size of the restricted refutation is at most the
degree/size of the original refutation.

We also consider more general substitutions.

\begin{definition}[Polynomial Substitution]
  Functions
  $\rho: \set{x_1, \ldots, x_n} \rightarrow \R[x]_{\le k}$ that map
  variables to polynomials of degree at most $k$ are called
  \emph{polynomial substitutions}.
\end{definition}

For polynomial substitutions we have the following well-known lemma.

\begin{lemma}
  \label{lem:sos degree polynomial substitution}
  Let $\calP$ be a system of polynomial equations and let $\rho$ be a
  polynomial substitution mapping variables to polynomials of degree
  at most $k$. Then,
  $\RefuteDeg{SoS}{\calP} \ge
  \RefuteDeg{SoS}{\restrict{\calP}{\rho}}/k$.
\end{lemma}

This lemma can again be verified by considering a refutation of
$\calP$. Substitute each variable $x_i$ in the proof by
$\rho(x_i)$. This results in a refutation of $\restrict{\calP}{\rho}$,
whose degree is at most a factor $k$ larger than the degree of the
refutation of $\calP$.

\subsection{The Circuit Size Formula}
\label{sec:circuit-encoding}
The formula $\circuit_s(f)$ encodes the claim that the function $f$,
given as a truthtable $f \in \set{0, 1}^{2^n}$, can be computed by a
circuit of size $s$ over $n$ Boolean inputs $x_1, \ldots, x_n$. The
encoding is not essential but for concreteness let us fix one encoding
of this claim. We deviate from the encoding used by Razborov
\cite{Razborov98, Razborov04ResolutionLowerBoundsPM} and do not
present the formula as a propositional formula but rather as a system
of polynomial equations. In order to encode below constraints as a
constant width CNF formula, as done by Razborov, one needs to
introduce extension variables. Despite this difference it is not
difficult to see that our lower bound also works against the CNF
encoding. In \cref{sec:encoding} we directly show that a low degree
SoS refutation of the CNF encoding gives rise to a low degree SoS
refutation of the encoding used in this paper (see
\cref{prop:cnf-poly}). Thus a lower bound for our encoding implies a
lower bound for the CNF encoding. As the presentation is simpler in
the polynomial encoding, we present it as follows.

We also need to define the monotone version of $\circuit_s(f)$ denoted
by $\circuit^{\mathsf{mon}}_s(f)$. The later is a restriction of the
former with the $\isneg(v)$ (see below) variable, for all $v \in [s]$,
set to $0$. This forces the circuit to only contain $\land$ and $\lor$
gates, i.e., the circuit is monotone.

All variables introduced in the following are Boolean variables and we
implicitly add the Boolean axiom $y(1 - y) = 0$ for each variable $y$
and further implicitly introduce the ``bar variable'' $\bar{y}$ along
with the negation axiom $y = 1 - \bar{y}$ (and the corresponding
Boolean axiom) ensuring that $\bar{y}$ is always the negation of $y$.

Let us first describe the \emph{structure variables} which are used to
describe the circuit that supposedly computes the function $f$.

We view the $s$ gates as being indexed from $1$ to $s$ in topological
order with gate $s$ being the output.
For each gate $v \in [s]$ there are three variables $\isneg(v),
\isor(v), \isand(v)$ indicating the operation computed at
$v$. Similarly for a gate $v \in [s]$ and a wire $a \in \set{1,2}$ we
have variables $\isfromconst(v, a), \isfromvar(v, a), \isfromgate(v, a)$
indicating whether the input wire $a$ of $v$ is connected to a
constant, a variable or a gate.

Further, we have the variables $\constval(v, a)$, $\isvar(v, a, i)$
and $\isgate(v, a, u)$, for $a \in \set{1,2}$, $i \in [n]$ and
$u < v$, specifying the constant value, input $x_i$ or gate $u$, the
input wire $a$ of $v$ is connected to (assuming $a$ is connected to
the corresponding kind).

The second set of variables are the \emph{evaluation variables}, which
describe what value is computed at each $v$ on input
$\alpha = \alpha_1, \ldots, \alpha_n$ (i.e., we have $x_i = \alpha_i$).

For each gate $v \in [s]$ and assignment $\alpha \in \set{0,1}^n$ we
have a Boolean variable $\outwire_\alpha(v)$ indicating the value
computed at gate $v$ on input $\alpha$. The Boolean variable
$\inwire_{\alpha}(v, a)$ indicates the value brought to the vertex
$v \in [s]$ on wire $a \in \set{1,2}$ on input $\alpha$.

Note that there is a total of
$3s + 6s + 2s + 2sn + 2{s \choose 2} = \Theta(s^2+sn)$ structure
variables, and a total of $3 s 2^n$ evaluation variables, for a total
of $\Theta(s^2 + s 2^n)$ variables in $\circuit_s(f)$.

The formula consists of the following axioms. For the sake of
readability we omit some universal quantifiers: the variable
$a \in \set{1,2}$ in
\cref{eq:ax-from,eq:ax-in-gt,eq:ax-in-gt2,eq:prod-zero-in,eq:prod-zero-gate,eq:inwire-const,eq:inwire-var,eq:inwire-gate}
as well as the variable $\alpha \in \set{0,1}^n$ in
\cref{eq:inwire-const,eq:inwire-var,eq:inwire-gate,eq:ax-neg,eq:ax-or,eq:ax-and,eq:ax-correct}
are implicitly universally quantified.

Let us first describe
the axioms on the structure of the circuit. In the following
section we refer to this set of axioms as the \emph{structure
  axioms}. The first axioms ensure that every wire is connected to a
single kind
\begin{align}
  \label{eq:ax-from}
  \isfromconst(v, a) +
  \isfromvar(v, a) +
  \isfromgate(v, a) = 1 \quad
  \forall \, v \in [s] \eqcomma
\end{align}
and similarly the next axioms make sure that each gate is of
precisely one kind
\begin{align}
  \label{eq:ax-fn}
  \isneg(v) + \isor(v) + \isand(v) = 1
  \quad \forall \, v \in [s] \eqperiod
\end{align}
The final structure axioms ensure that the variables, which indicate
to what input or gate a fixed wire is connected to, always sum to one
(except for gate $1$ which cannot have any inputs from other gates)
\begin{align}
  \label{eq:ax-in-gt}
  \sum_{i = 1}^{n}\isvar(v,a,i) &= 1
  \quad \forall v \in [s], \text{ and}\\
  \label{eq:ax-in-gt2}
  \sum_{u = 1}^{v-1}\isgate(v,a,u) &= 1
  \quad \forall v \in [s] \setminus\{1\} \eqperiod
\end{align}
We further strengthen our encoding by adding the axioms
\begin{align}
  \label{eq:prod-zero-in}
  \isvar(v,a,i)\isvar(v,a,j) &= 0
  \quad\forall v \in [s],\, i < j \in [n], \text{ and}\\
  \label{eq:prod-zero-gate}
  \isgate(v,a,u)\isgate(v,a,u') &= 0
  \quad\forall u < u' < v \in [s] \eqperiod
\end{align}
Note that \cref{eq:prod-zero-in,eq:prod-zero-gate} are implied by
\cref{eq:ax-in-gt,eq:ax-in-gt2}. We add these axioms in order to argue
that a short refutation of the CNF encoding of this principle leads to
a short refutation of the present encoding.

The second group of axioms are the \emph{evaluation axioms} and they
ensure that the evaluation variables indeed compute the intended
values. We start by making sure that the wires carry the value
intended by the structure axioms. If a wire is connected to a
constant, then the evaluation variable associated with that wire
should always be equal to the constant
\begin{align}
  \label{eq:inwire-const}
  \isfromconst(v,a) \cdot
  \big( \inwire_{\alpha}(v,a) - \constval(v, a) \big) = 0 \eqcomma
\end{align}
and similarly in case if a wire is connected to an input or a gate
\begin{align}
  \label{eq:inwire-var}
  \isfromvar(v, a) \cdot \isvar(v, a, i) \cdot
  \big( \inwire_\alpha(v, a) - \alpha_i\big) &= 0
  \eqcomma\\
  \label{eq:inwire-gate}
  \isfromgate(v, a) \cdot \isgate(v, a, u) \cdot
  \big(\inwire_\alpha(v, a) - \outwire_\alpha(u)\big) &= 0 \eqperiod
\end{align}
The final set of evaluation axioms makes sure that the output
evaluation variable of a gate is correctly related to the input
evaluation variables:
\begin{align}
  \label{eq:ax-neg}
  \isneg(v) \cdot
  \outwire_{\alpha}(v)
  &=
    \isneg(v) \cdot
    \overline{\inwire_{\alpha}(v, 1)}\eqcomma\\
  \label{eq:ax-or}
  \isor(v) \cdot
  \outwire_{\alpha}(v)
  &=
    \isor(v)
    \cdot \big(
    1 -
    \overline{\inwire_{\alpha}(v, 1)} \cdot
    \overline{\inwire_{\alpha}(v, 2)}
    \big) \eqcomma\\
  \label{eq:ax-and}
  \isand(v) \cdot
  \outwire_{\alpha}(v)
  &=
    \isand(v) \cdot
    \inwire_{\alpha}(v, 1) \cdot
    \inwire_{\alpha}(v, 2) \eqperiod
\end{align}

Last but not least we have the axioms that ensure that the circuit
outputs the function specified by the truthtable
\begin{align}
  \label{eq:ax-correct}
  \outwire_{\alpha}(s) &= f(\alpha) \eqperiod
\end{align}
